% 第1章正文，根据《一二章.md》更新。
\FloatBarrier\clearpage

\section{作品概述}\label{chapter:1}

\subsection{研究背景与作品目标}\label{section:1.1}

\subsubsection{网络信息环境下的个人认知安全问题}\label{section:1.1.1}

互联网已成为人们了解社会、获取知识和作出判断的重要渠道。无论是学生选择升学方向、求职者了解就业形势，还是普通用户关注社会热点，人们都越来越依赖社交平台和短视频平台获取信息。这些信息不仅影响人们对客观事实的了解，也为个人形成判断和作出决策提供依据。因此，个人能否接触到充分且可靠的信息，直接关系到其认识外部世界和处理现实问题的能力。

随着网络信息规模不断扩大，个性化推荐逐渐成为互联网平台组织和分发内容的重要方式。推荐系统通过分析用户的浏览记录和互动行为，向其推送更符合兴趣的内容，在提高信息获取效率的同时，也可能使用户反复接触某些特定类型的信息。与此同时，人们通常更愿意关注自己熟悉或认同的内容，并主动选择相关账号和信息来源。平台推荐与用户主动选择相互作用，可能使个人看到的信息逐渐集中于少数主题或立场。当用户长期依赖这些内容了解外部世界时，便可能忽视其他观点和事实，从而形成相对片面的认识。

这种现象值得关注，并不是因为持续阅读同类内容本身存在问题，而是因为平台上某类事件反复出现，并不意味着它在现实中同样普遍。例如，一名求职者如果长期看到关于企业裁员和毕业生求职失败的报道，却很少接触不同行业的招聘情况，便可能高估就业困难的普遍程度，将部分群体的经历理解为整体就业市场的状况。特别是当大量内容反复引用同一则新闻时，用户还可能误以为这些报道提供了相互独立的事实支持，进而高估某一判断的可靠性。

当个人长期依赖范围有限的信息作出判断时，便可能忽视问题的其他影响因素，使其对现实情况的认识不够全面。这种认识偏差在就业和升学等需要综合信息进行决策的场景中，可能影响个人对风险与机会的评估，甚至导致不必要的担忧或不合理的决策。

当然，长期关注特定主题并不一定会产生认知偏差。用户也可能通过集中阅读获得充分而深入的专业知识。因此，个人认知安全保护的重点，不应是简单限制用户浏览某类内容，而是帮助用户判断自己看到的信息是否过于单一，某种观点的出现频率是否能够反映现实情况，以及是否忽略了其他值得了解的信息。

\subsubsection{信息茧房的形成及其个人识别难题}\label{section:1.1.2}

信息茧房通常用于描述个人长期处于相对单一的信息环境，频繁接触符合既有兴趣或立场的内容，而较少了解其他观点和信息的现象。它既可能受到个性化推荐机制的影响，也可能源于用户主动选择信息或长期关注特定账号等行为。推荐算法会利用用户过去的互动记录调整推送内容，用户又可能倾向于选择与自身认识一致的信息，两者相互作用，使个人接触不同观点的机会逐渐减少。然而，这一变化往往发生在日常浏览过程中，普通用户难以及时察觉，主要有以下三个原因。

\textbf{（1）相同信息可能以不同形式反复出现，使用户难以察觉信息的重复传播。} 同一事件可能被不同账号反复传播，并被包装成标题和表达方式各异的内容。例如，用户连续观看十条关于就业困难的视频，虽然这些视频来自不同博主，但其中多条可能引用同一则裁员新闻。如果用户无法辨认它们之间的联系，就容易将同一事实的重复传播误认为多个独立证据，进而高估某种观点的可靠性，而忽略其他尚未接触的事实和解释。

\textbf{（2）浏览内容分散在不同时间，用户难以发现信息长期重复的现象。} 短视频和帖子通常以连续推荐、快速切换的方式呈现，同一主题的内容可能分散在不同时间，中间穿插大量无关信息。这种浏览方式带来两方面的困难：其一，用户主要关注当前内容，难以将其与过去看到的信息持续比较。当不同账号长期重复使用相同材料时，用户可能意识不到这种重复，更难判断这一现象持续了多长时间；其二，用户难以将分散出现的观点和论据联系起来。即使阅读了大量内容，也未必能够准确回忆某个问题上出现过哪些不同观点，以及各自获得了哪些事实支持，最终可能只记住最熟悉或情绪最强烈的说法。

\textbf{（3）浏览主题集中不等于信息茧房，用户难以区分正常阅读与信息单一的情况。} 例如，两名用户都连续数日关注就业问题，其中一人不断了解新的招聘数据和不同行业的岗位情况，另一人却反复观看引用相同裁员新闻的负面报道。尽管两人的浏览主题相同，但前者持续获得新的信息，后者接触的事实依据则较为有限。如果只统计浏览主题和内容数量，就难以区分正常的专题阅读与长期接触重复信息的情况。

因此，要帮助个人识别不易察觉的信息茧房，不能仅依靠内容数量和主题分布的统计，还需要进一步辨认不同内容表达的观点及其论据，并与历史浏览记录进行比较，判断信息重复和观点集中是否长期持续。同时，还应将零散内容整理为能够回顾和追溯的观点总结，帮助用户了解自己实际获得了哪些信息、可能遗漏了哪些判断依据。

\subsubsection{作品总体目标}\label{section:1.1.3}

针对个人难以察觉信息茧房、难以整理零散浏览内容，以及反复接触过度推断型情绪引导内容可能带来的认知风险，本作品研发了面向个人认知安全的智能体检工具``破茧''。作品以用户日常浏览的网络内容为分析对象，通过人工智能技术识别信息茧房风险及过度推断型情绪引导的接触情况，并利用可验证机密计算技术保护用户隐私，为个人提供安全可信的认知体检服务。

在信息茧房识别方面，破茧不仅关注用户每天浏览了多少内容、花费了多少时间，还进一步分析用户在同一主题下看到了哪些不同观点，以及这些观点是否反复使用相同的论据。通过对长期浏览记录的分析，系统能够区分持续获得新信息的集中阅读与长期重复接触相似内容的情况，识别用户是否长期局限于少数观点，且缺少新的事实和解释。通过系统对用户浏览内容的持续分析，帮助用户及时发现自身信息获取的局限，避免因长期接触片面信息而影响对现实问题的认识。

对于用户已经阅读却难以记住的大量零散内容，破茧进一步提供主题观点整理功能。系统能够按照用户关注的主题，归纳当天接触到的主要观点，并统计每种观点获得了多少条已浏览内容的支持。同时，整理相关观点采用的主要论据，使用户无需重新翻看大量视频和文章，也能够回顾自己已经获得的信息。

作品还将认知体检从被动的结果展示延伸到主动提醒。当系统发现用户长期浏览重复信息或反复接触以局部事实夸大负面结论的内容时，可以结合具体浏览记录提供提示。对于观点明显集中或缺少其他论据的主题，系统还能够指出已浏览信息中存在的局限，为用户主动补充资料提供参考。

为保护个人浏览数据的隐私安全，破茧在持续认知分析过程中引入可验证机密计算技术。上述信息茧房识别、情绪引导识别和主题观点整理功能需要长期处理用户的浏览记录，而这些记录可能反映其兴趣偏好和现实处境，属于敏感用户隐私数据。为减少敏感信息在处理过程中的暴露，系统引入可验证机密计算技术，将浏览数据的解密与分析限制在经过验证的可信执行环境中，并由客户端在验证计算环境及程序身份后提交加密数据。在真实硬件隔离、远程证明及密钥管理机制有效的条件下，包括作品运营团队在内的服务提供者也无法仅凭普通服务器管理权限读取受保护的明文浏览记录。

最终，破茧通过小时报告、每日报告和长期分析报告，将信息茧房识别与情绪引导识别结果转化为可以查看的个人认知体检信息，并提供相应的原始内容作为判断依据。作品希望帮助用户看清自己实际接触的信息，整理已经阅读的不同观点，发现可能存在的信息局限，并在保护个人隐私的前提下更加主动地管理自身的信息环境。

\subsection{作品拟解决的关键问题}\label{section:1.2}

围绕个人信息茧房识别与认知安全保护，本作品重点解决三个相互关联的问题：如何区分正常的信息浏览与可能存在的信息茧房，如何识别网络内容中通过不合理推断强化负面情绪的表达，以及如何在持续处理个人浏览记录的过程中保护用户隐私。

\subsubsection{如何区分正常的信息浏览与可能存在的信息茧房？}\label{section:1.2.1}

信息茧房识别的核心难点，在于如何区分正常的集中浏览与长期接触单一信息的情况。用户频繁关注某个话题，并不意味着其已经陷入信息茧房。例如，两名用户都连续多天浏览就业相关内容，其中一人不断获取新的招聘数据，了解不同地区的就业情况；另一人虽然浏览了大量视频，却反复看到相同的观点和新闻事件。如果仅统计浏览主题和观看数量，两人的行为可能十分相似，但实际获得的信息却存在明显差异。因此，系统不能简单地将浏览主题集中等同于信息茧房，还需要考察同一议题下的观点是否趋于单一，以及支撑这些观点的论据是否缺乏增量。

要比较用户看到的观点与论据，首先需要从浏览内容中还原具体的事实与判断。网页正文、视频字幕和口播可能分别提供信息，同一段内容还可能同时引用他人观点并表达自己的立场。针对这一问题，本作品采用\textbf{基于轻量化信息抽取的观点增量识别方法}：先将网页文字、OCR识别结果和语音转写文本按原有顺序整理，再使用UIE-base定位议题、观点、论据及其适用条件；随后由基于Qwen3.5-0.8B的Simple Jev判断语句类型，以及观点与论据的支持或反驳关系。由此形成带有原文位置和来源的结构化记录，为个人历史比较提供统一输入。

获得单条内容的观点与论据后，还需要确定这些材料相对于个人历史新增了什么。同一观点可能被不同事实支持，而不同账号也可能将同一事实换一种说法反复传播。为此，破茧先利用BM25与BGE-small-zh-v1.5召回相近议题，并由基于Qwen3.5-4B的SemIf判定议题归属；随后在所属议题下分别检索相近观点和历史论据，判断观点能否归并、事实是否等价以及适用条件是否一致。对于一条论据包含的多个事实要点，系统综合不同历史记录的覆盖情况计算信息新颖度，从而识别真正新增的判断依据。

在识别每次浏览的信息增量后，还需要判断观点集中与依据重复是否持续存在。用户短期围绕某一任务阅读大量类似内容，并不必然形成信息茧房。为此，系统按议题和时间窗口计算观点集中度、有效观点数量、去重论据数量及信息增量率，并考察主导观点持续集中且新增依据持续不足的情况。风险判定采用连续有效窗口的联合条件，输出关联原始来源的预警结果；对于证据不足的窗口，则保留尚不能判定的状态。

上述方法形成从内容结构化抽取、同议题历史关系判别到长期信息增量统计的连续技术链。每一步保留原文依据并将结果交给下一阶段，既支持信息茧房预警，也可直接用于主题观点整理。通过对固定候选类别进行快速评分、缩小历史检索范围并复用重复内容的处理结果，系统减少高频浏览场景中的重复计算。

\subsubsection{如何识别网络内容中的情绪引导及其潜在认知风险？}\label{section:1.2.2}

网络信息不仅通过事实和观点影响人们对问题的认识，其表达方式也可能影响人们关注哪些风险。在就业和升学等与个人生活密切相关的话题中，一些内容通过数据介绍现实困难让人们感到焦虑，另一些内容则可能借助夸张标题和极端案例激发强烈情绪。由于情绪化内容容易吸引关注并引发互动，部分内容可能在用户主动传播和推荐算法的共同作用下获得更多曝光，使用户反复看到带有相似情绪的表达。

这种传播过程可能使用户长期接触某一类情绪化内容，进而影响其对相关问题的认识。例如，两条视频都在讨论就业困难，一条主要介绍不同行业的招聘数据，另一条则反复强调``毕业就失业''\,``普通人没有出路''等判断。后者通过绝对化语言将特定困难描述为普遍且难以改变的处境。如果用户长期反复接触类似内容，便可能更加关注求职失败等负面结果，而忽视其他行业和地区的实际情况，进而对自身就业前景产生不必要的担忧。

识别此类内容，关键不在于是否出现``焦虑''或``恐惧''等词语，而在于它如何从已知事实推导出负面结论。正常的风险提醒也可能包含负面信息，但结论仍符合事实的适用范围；具有引导性的表达则可能把局部事件推向普遍判断，或将尚不确定的风险表述为必然结果。因此，系统需要同时考察事实与结论之间的推断关系，以及相关表达是否进一步强化负面情绪。

针对上述问题，本作品提出\textbf{基于共享模型多专家判别的情绪引导识别方法}。系统复用前一阶段提取的事实、观点及论证关系，以整条内容为上下文，对其中每项评价性结论分别进行判断。基于Qwen3.5-9B设置的三个语义专家共享模型及输入上下文，通过Simple Jev对预设类别快速评分，分别识别事实适用范围是否被扩大、结论确定程度是否被提高，以及是否存在灾难化或无助化等负面情绪强化。随后根据联合规则判断是否构成过度推断型情绪引导，并结合原文位置与语境核验相关依据。

单条结论的识别结果还需要与用户实际浏览行为关联，才能用于长期认知体检。系统将结论级判别汇总至浏览事件，再结合所属议题、接触时间和有效浏览时长，统计情绪引导内容的出现比例。若用户长期浏览将个别求职困难推及整体就业前景的内容，系统既可以展示此类内容的接触频率，也可以提供相关的事实片段与结论原文，说明具体的外推方式。上述结果与观点分布及论据增量共同构成个人认知体检信息。

这些指标描述的是浏览内容的推断方式、负面情绪强化特征及用户接触情况，不直接衡量用户实际产生的情绪或心理变化。破茧通过具体的原文依据展示识别结果，为用户判断相关内容提供参考。

\subsubsection{如何在保护用户隐私的前提下完成持续的认知分析？}\label{section:1.2.3}

个人认知体检需要分析较长时间内的浏览记录，这使其面临比普通内容分类任务更加突出的隐私保护问题。单条浏览记录可能只反映用户对某个问题的短暂关注，但当大量记录按照时间顺序积累后，就可能揭示用户长期关注的事项和生活处境。例如，连续浏览求职或债务相关内容，可能使外部观察者推测用户正在面临的现实问题。即使单条记录并不敏感，长期关联分析也可能形成具有较高隐私价值的个人画像。

如果这些数据被未经授权的人员获取，便可能用于分析个人偏好或开展针对性营销，甚至对用户实施欺诈。因此，破茧在开展长期认知分析的同时，必须避免将用户原本分散在不同平台上的私人浏览记录集中暴露给新的数据处理方。特别是作为认知体检服务的运营者，系统开发团队本身也不应拥有随意读取用户浏览内容的能力，否则用户仍需要完全信任服务提供者的管理行为，难以从根本上解决隐私顾虑。

传统的网络服务通常通过传输加密和服务器访问控制保护用户数据。然而，复杂的人工智能模型往往需要在服务器上处理明文输入，即使数据在传输和存储过程中经过加密，服务器在执行模型推理时仍可能接触解密后的内容。这意味着，拥有相应权限的服务器管理人员可能通过访问程序内存或修改数据处理流程获取敏感信息。因此，仅依靠服务器运营方承诺不查看数据，并不能从技术上保证用户隐私不被服务提供者获取。

针对敏感信息在不同环节可能发生的泄露，破茧构建了\textbf{基于可验证异构机密计算的全过程隐私保护方法}，围绕手机采集、数据上传、CPU与GPU计算及历史存储设置防护。客户端通过公开安全关键代码、可复现构建和经批准的发布版本约束采集与对外通信行为。云端则将解密及模型推理限定于经过验证的可信工作负载。``异构''是指CPU和GPU均处于受保护的计算范围，避免只保护服务器内存，却在显卡推理阶段暴露用户数据。

在数据上传前，手机客户端首先核验云端Intel TDX可信虚拟机与NVIDIA机密GPU的远程证明，并检查实际运行状态是否符合已批准的程序版本。可信程序在受保护内存中生成临时密钥，客户端核验该公钥与本次证明的绑定关系，验证通过后才使用HPKE加密提交浏览内容。普通网关只转发密文，真正的解密由可信程序在受保护的CPU内存中执行。此后的OCR、语音转写与语义模型推理均在受保护的CPU和GPU范围内完成。

\textbf{由此，包括作品运营团队在内的服务提供者也不能仅凭普通服务器管理权限获取用户浏览明文。} 这一限制既依赖CPU与GPU的硬件隔离，也依赖经过审核的程序版本和受约束的通信行为。运营方即使控制普通业务网关，也不能通过自行替换程序或接收公钥，在客户端验证未通过的情况下取得解密权限。

单次计算结束后，用户的历史记录及分析结果还需受到持续保护。破茧将长期主密钥保留在用户手机，并使用Android Keystore提供本地保护；历史索引、语义记录和认知体检报告均通过AES-256-GCM加密存储。每次进行历史比较前，手机重新核验可信程序，并按次交付本次所需的数据密钥。可信程序在受保护环境内读取、更新并加密写回历史状态，客户端通过保存的版本标识核对更新连续性。由此，普通数据库可以管理密文，却不能凭存储权限恢复长期分析记录。

通过客户端程序审核、云端远程证明、CPU与GPU受保护推理及用户控制的历史密钥，破茧使隐私保护覆盖连续认知分析的主要环节。这一保护建立在可信硬件、程序审核机制与密钥隔离有效的条件下，使普通服务器运维权限不能直接转化为个人浏览数据的读取权限。

上述三个问题共同构成破茧的研究基础。作品以观点集中与论据增量识别信息茧房，以事实和结论的关系判别识别过度推断型情绪引导，并以可验证异构机密计算保护连续分析涉及的敏感数据。内容抽取结果既用于个人历史比较，也为情绪引导识别和观点整理提供统一依据，最终形成能够追溯原始内容的个人认知体检报告。

\subsection{国内外研究现状及现有方案不足}\label{section:1.3}

\subsubsection{信息茧房的形成机制与识别研究}\label{section:1.3.1}

信息茧房是国内外传播学和计算社会科学领域长期关注的问题。相关研究主要从平台推荐机制、用户信息选择和社交网络结构等角度展开，分析为什么用户容易接触与自身立场相近的内容，以及这种现象是否会影响个人对现实问题的认识。

在信息茧房的形成机制方面，Bakshy 等于 2015 年在《Science》发表研究，通过分析超过一千万名 Facebook 用户的数据，考察社交关系和用户主动选择对信息获取范围的影响\cite{ref3}。研究发现，社交关系和用户自身的点击选择都可能限制不同政治立场的信息接触，其中主动选择具有重要作用。这说明信息茧房并非完全由推荐算法造成，即使平台提供了不同观点的内容，用户也可能主动选择符合自身立场的信息。

Robertson 等于 2023 年在《Nature》发表研究，将用户在 Google 搜索中实际看到的新闻结果与最终点击的内容进行比较\cite{ref4}。研究发现，用户主动点击的新闻比搜索结果中展示的新闻更符合其政治立场。这进一步说明，平台向用户展示了什么，与用户最终选择阅读什么，并不是同一个问题。因此，仅分析推荐结果或平台传播结构，仍难以完整反映个人实际接触的信息。

在信息茧房的量化识别方面，Cinelli 等于 2021 年在《PNAS》发表研究，通过比较多个社交平台的用户互动关系与信息传播情况，揭示了不同平台中相近观点聚集和传播的特点\cite{ref1}。这类研究通常通过用户之间的连接关系和传播网络中的群体聚集程度，判断是否存在观点相近的用户群体。相关方法能够揭示平台整体的信息传播特征，为理解信息茧房的形成提供重要依据。

国内学者也对信息茧房的形成机制和识别标准进行了深入讨论。陈昌凤和仇筠茜指出，信息茧房的概念在实际应用中存在泛化和误读的问题，不能将算法推荐与信息茧房直接画等号\cite{ref34}。虞鑫和王金鹏进一步研究智媒时代的信息茧房问题，讨论技术机制与用户行为之间的关系\cite{ref6}。这些研究说明，信息茧房的形成具有复杂性，不能简单依据平台算法特征或使用时间作出判断。

此外，信息接触范围与认知结果之间也并不存在必然的对应关系。Nyhan 等于 2023 年在《Nature》发表随机干预研究，通过减少 Facebook 用户接触立场相近来源的内容，观察其政治态度和相关行为变化\cite{ref5}。研究发现，在实验观察期内，信息来源的调整并未使所测量的政治态度和极化指标发生显著变化。这表明，信息环境的单一性与实际认知影响需要分别研究，不能仅凭用户接触相似观点就断言其已经形成认知偏差。

总体来看，已有研究能够从平台和群体层面分析观点聚集现象，也揭示了用户主动选择的重要作用，但将这些成果直接用于个人信息茧房识别仍存在困难。一方面，用户关注特定主题可能是正常的学习或信息查询行为，主题集中并不意味着信息缺乏多样性。另一方面，不同账号可能反复传播同一事件，即使信息来源看起来较为丰富，用户实际获得的判断依据仍可能十分有限。

因此，面向个人的信息茧房识别，不能只统计浏览主题或信息来源，还需要分析具体内容表达的观点及其论据，并进一步判断这种信息单一现象是否长期持续。这也是本作品开展个人浏览内容分析的主要研究切入点。

\subsubsection{多模态观点理解与历史信息比较研究}\label{section:1.3.2}

准确识别个人信息茧房，首先需要理解用户实际浏览的内容。近年来，自然语言处理领域围绕观点识别、论证分析和多模态语义理解开展了大量研究，为从复杂网络内容中提取观点及其论据提供了技术基础。

在观点识别方面，Mohammad 等提出的 SemEval-2016 立场检测任务，要求模型判断文本对特定目标持支持或反对等立场\cite{ref8}。这一研究推动了面向具体问题的观点识别方法发展，使系统能够区分内容讨论的主题与实际表达的判断。随后，Liang 等于 2022 年提出 JointCL 联合对比学习框架，通过立场对比学习与目标感知的原型图对比学习，提高模型对未见过的讨论目标的泛化能力\cite{ref15}。

随着社交平台内容形式不断丰富，研究者开始将观点识别扩展到图文联合分析。Zhang 等于 2025 年在 EMNLP 发表 T-MAD 模型，通过目标驱动的跨模态对齐与动态权重调整，增强模型对文本和图像之间关系的理解能力\cite{ref14}。该方法能够根据讨论目标融合不同模态的信息，缓解文本与图像表达不一致所带来的判断困难。

然而，仅识别内容支持或反对某一观点，仍不足以判断用户实际获得了多少不同的信息。例如，两篇文章即使都认为就业形势严峻，也可能分别引用招聘统计和企业经营数据。若模型只输出相同的立场标签，就会忽略两篇文章在论据上的差异。

针对这一问题，论证挖掘研究进一步关注观点与论据之间的关系。Cheng 等于 2022 年在 ACL 发表 IAM 数据集及相关研究，将观点提取和证据提取等任务纳入统一的论证分析框架\cite{ref13}。该研究为自动识别文章中的判断及其论据提供了数据和方法基础，使内容分析不再局限于简单的主题分类。

在此基础上，如何区分新增信息与已有信息的重复，也成为信息检索领域的重要研究方向。美国国家标准与技术研究院组织的 TREC 新颖性检测研究，要求系统在连续出现的相关文档中找出此前没有出现过的信息\cite{ref35}。这类研究说明，信息是否具有新意，需要结合已经阅读过的内容进行判断，而不能只分析当前文本本身。

不过，将上述方法应用于个人浏览历史时，还需要解决更加复杂的信息比较问题。不同账号可能使用不同措辞引用同一则新闻，也可能在表达相同观点的同时提供新的统计数据。单纯依赖文本相似度，容易将措辞变化误认为新增信息，也可能忽略观点相同但论据不同的情况。因此，需要同时比较内容的观点和论据，明确当前内容究竟重复了什么，又提供了哪些新的事实或解释。

此外，个人浏览记录会随着时间不断增加，信息之间的关系也会持续变化。Rossi 等提出的 Temporal Graph Networks 通过记忆模块与图结构学习，对连续发生的动态事件进行建模\cite{ref16}。这表明连续浏览历史可以借助持续更新的状态开展分析。对个人信息茧房识别而言，还需要直接衡量同一议题下的观点分布与事实依据增量，才能将时间变化转化为明确的风险判定。

综合现有研究可以发现，观点识别和历史信息比较已经形成了重要的技术基础，但这些方法通常针对各自独立的分析任务。面向个人信息茧房识别，需要进一步将不同形式的内容理解结果统一起来，比较其与历史材料的关系，并分析信息重复现象是否长期持续。

此外，个人认知体检还需要将分析结果转化为用户容易理解的总结。例如，用户可能希望知道自己当天围绕某个主题看到了多少种观点，每种观点有多少条内容支持，以及这些内容使用了哪些论据。这不仅要求系统准确识别观点，还需要将分散在不同内容中的相近判断合并整理，并保留对应的原始来源。

因此，本作品将多来源浏览文本的观点论据提取与个人历史信息比较相结合，不仅用于识别信息重复，还进一步支持主题观点整理和分析依据追溯，使相关研究成果能够服务于个人日常信息管理。

\subsubsection{情绪化内容传播与细粒度情绪理解研究}\label{section:1.3.3}

除了信息内容本身，网络信息的情绪表达也受到国内外研究者的广泛关注。相关研究主要围绕情绪对信息传播的影响，以及如何通过人工智能技术识别情绪及其产生原因展开。

在情绪传播方面，Berger 和 Milkman 于 2012 年在《Journal of Marketing Research》发表研究，通过分析《纽约时报》文章的传播情况，发现能够激发强烈情绪的内容通常更容易获得分享\cite{ref36}。其中，愤怒和焦虑等高唤醒情绪与较强的传播能力有关，而悲伤等低唤醒情绪并不表现出相同的规律。这说明信息传播不仅受到内容主题的影响，也与内容能否引发用户情绪反应密切相关。

Brady 等于 2017 年在《PNAS》发表研究，进一步发现，在特定政治与道德议题中，包含道德情绪表达的信息更容易在社交网络中传播，而且这种传播作用在立场相近的群体内部更加明显\cite{ref9}。这一结果表明，情绪表达可能与群体内部的信息反复传播相互联系，使某些观点在特定用户群体中持续获得关注。

推荐算法也可能进一步放大这一现象。Milli 等于 2025 年在《PNAS Nexus》发表算法审计研究，通过比较基于用户互动的排序机制与按时间排序的内容展示方式，发现前者更容易放大带有愤怒和群体对立情绪的内容\cite{ref37}。这一结果说明，在特定平台环境下，互动量较高的情绪化内容可能获得额外曝光，而高互动量并不必然意味着用户更喜欢这些内容。

这些研究揭示了情绪化内容容易传播并可能被推荐机制放大的原因。对于长期使用社交平台的个人而言，如果某类负面信息不断获得曝光，用户就可能反复接触相似的情绪表达，并对相关风险产生持续关注。因此，除了了解浏览内容讨论了什么问题，还需要分析这些内容具有怎样的情绪特点。

在情绪分析方法方面，早期研究主要通过情绪词典和文本分类模型识别内容中的情绪类别。然而，网络情绪并不总是直接通过情绪词语表达，单纯依赖情绪标签难以判断某种情绪由什么事件或判断引起。

针对这一问题，Ding 等于 2020 年提出 ECPE-2D 模型，通过二维联合表示与交互建模，直接提取文本中的情绪及其原因\cite{ref23}。Troiano 等于 2023 年进一步将认知评价理论引入文本情绪分析，研究事件评价与情绪表达之间的联系\cite{ref26}。这类方法使情绪识别从简单判断内容是否积极或消极，发展到进一步分析情绪与具体事件之间的关系。

随着多模态学习技术的发展，情绪原因分析也逐渐从纯文本扩展到视频和语音场景。SemEval-2024 Task 3 将多模态情绪原因分析作为研究任务，要求模型结合不同模态的信息识别情绪及其对应原因\cite{ref28}。相关研究为理解视频中的情绪表达提供了重要技术支持。

现有研究已经能够解释部分情绪化内容为何更容易传播，也形成了识别情绪及其原因的技术方法。但传播研究主要关注内容扩散规律，情绪理解研究则通常以单条文本或对话为分析对象。对于借助局部事实推导出绝对化负面结论的情绪引导，仅识别情绪类别及原因仍不足以说明问题，还需要比较事实与结论是否相符，并追溯其原文依据。

因此，本作品将事实适用范围、结论确定程度和负面情绪强化的联合判别与个人浏览历史结合起来，识别过度推断型情绪引导，并按议题统计相关内容的接触频率与浏览时长，为长期提醒提供可追溯的依据。

\subsubsection{面向个人数据分析的隐私保护研究}\label{section:1.3.4}

个人认知体检需要长期保存和分析浏览记录，因此，隐私保护不仅关系到数据安全，也直接影响用户是否愿意使用相关工具。尤其是在复杂人工智能模型依赖服务器计算的情况下，如何避免敏感数据在处理过程中泄露，已成为隐私计算领域的重要研究方向。

传统数据保护方法主要依靠传输加密和存储加密，减少数据在网络通信及持久化存储过程中的泄露风险。然而，服务器运行人工智能模型时，通常需要将加密数据恢复为可计算的形式。如果明文数据直接暴露于普通服务器内存中，就可能受到高权限软件或被入侵系统组件的影响。因此，仅依靠传统加密方式，难以完整解决敏感数据在计算过程中的保护问题。

针对这一问题，机密计算技术通过可信执行环境保护正在处理的数据。国际机密计算联盟将运行中数据的保护作为机密计算的核心目标，推动相关硬件隔离技术与软件框架的发展\cite{ref10}。可信执行环境能够为敏感程序提供受保护的运行空间，使其在符合硬件安全假设的条件下，减少外部管理组件直接访问明文数据的机会。

然而，可信执行环境本身并不足以解决用户信任问题。即使服务器支持机密计算，用户仍需要确认自己的数据是否真正交给了预期的程序，以及服务器运行环境是否符合安全要求。

针对这一需求，远程证明技术为客户端验证服务器的运行状态提供了方法基础。IETF 于 2023 年发布的 RFC 9334 系统规定了远程证明的基本架构，明确了证明方与验证方之间的职责，以及证明结果如何用于建立信任\cite{ref11}。通过对运行环境产生的可信证据进行验证，客户端能够在提交敏感数据前判断目标环境是否满足预期条件，从而减少单纯依赖服务提供方承诺所带来的信任风险。

随着大模型应用的发展，机密计算也开始与人工智能推理服务相结合。NVIDIA 发布的机密计算部署参考架构，将可信虚拟机与远程证明结合起来，为受保护环境中的模型运行提供了工程参考\cite{ref17}。相关方案说明，机密计算不仅可以用于传统的数据处理任务，也能够支持具有较高计算需求的人工智能应用。

不过，将机密计算应用于个人认知体检仍面临具体挑战。首先，系统需要持续处理新的浏览记录，并反复使用历史分析结果，因而必须保护长期保存的数据及其更新过程。其次，多模态模型通常涉及较大的计算开销，需要明确推理程序及相关数据是否处于可信保护范围之内。此外，远程证明只能验证规定范围内的运行环境，不能代替对应用程序安全性和数据使用行为的全面检查。

因此，本作品将客户端代码审核与版本验证纳入隐私保护链路，利用远程证明核验云端可信程序及CPU、GPU机密计算状态，并由手机管理长期密钥。云端在受保护环境内完成内容分析与历史比较，普通服务器仅转发或保存密文，使持续认知分析不再单纯依赖服务运营方的隐私承诺。

综上所述，国内外相关研究已经为信息茧房分析和情绪内容理解提供了重要基础，隐私计算技术也为保护个人敏感数据提供了新的实现条件。但面向个人日常浏览场景，仍需要将内容理解、历史比较和长期分析结合起来，并进一步解决分析结果难以解释以及个人浏览数据需要持续保护的问题。

针对这些需求，破茧以个人实际浏览记录为基础，在识别信息茧房风险的同时提供主题观点整理和情绪引导识别服务，并通过可验证异构机密计算技术保护敏感数据，形成兼顾认知安全与个人隐私的智能体检方案。

\subsection{作品主要功能与特色}\label{section:1.4}

\subsubsection{作品主要功能}\label{section:1.4.1}

破茧围绕用户日常浏览过程中存在的信息分散、内容重复和过度推断型情绪引导反复出现等问题，提供从个人信息体检到风险提醒的连续服务。同时，作品能够将已经阅读的内容整理成结构清晰的主题总结，并通过隐私保护机制减少个人浏览记录的暴露风险。

\textbf{（1）个人信息体检：浏览时间分布与信息获取特征分析}

在日常使用网络时，用户可能花费数小时浏览短视频和社交平台，却很难在一天结束后准确回忆自己看到了哪些重要信息。大量零散内容还可能造成一种错觉：虽然浏览了许多视频，实际上获得的新信息却十分有限。

针对这一问题，破茧提供个人信息体检功能。系统对用户授权范围内的浏览记录进行分析，通过24小时浏览分布图展示不同时段的主要关注主题，并统计不同观点的出现情况。同时，将当前内容与历史记录进行比较，帮助用户区分重复出现的信息与新获得的事实依据。

例如，用户在一天内观看了30条有关就业形势的视频。破茧不仅能够展示这些视频占用了多少时间，还能够分析其中有多少内容反复引用相同的新闻事件，以及哪些视频提供了新的招聘数据。用户由此能够更加直观地了解自己一天的信息获取情况，而不再只能通过观看数量或使用时间评价浏览收获。

\textbf{（2）信息茧房预警：长期观点集中与论据重复风险识别}

信息茧房最难察觉的地方，在于用户可能每天都在看到新视频，却没有意识到这些内容实际上反复表达相同的观点。当不同账号使用不同标题和表达方式传播相同材料时，用户还可能误以为某种说法得到了大量独立证据的支持。

针对这一问题，破茧提供信息茧房识别与主动预警功能。系统持续比较个人浏览历史，分析用户关注的观点是否过于集中，以及不同内容是否长期重复使用相同论据。当某个主题持续出现观点单一且信息重复的情况时，系统向用户提供提醒，并说明相关判断的具体依据。

例如，当用户连续多天观看有关就业困难的视频，而不同博主反复引用的都是同一则裁员新闻时，系统可以提示近期浏览内容主要集中于少数观点，且多条视频重复使用了相同材料。用户还可以查看具体内容，了解系统为什么作出这一判断。

对于正常的专题学习或集中信息查询，系统则结合是否出现新的事实依据及浏览情况的长期变化进行区分，避免仅凭某个主题出现频繁就发出预警。通过这一功能，用户能够及时发现个人信息环境中值得关注的变化，减少因重复浏览而高估某种观点可靠性的风险。

\textbf{（3）情绪引导识别与提醒：过度推断及负面强化的证据化展示}

在短视频和社交平台中，能够激发焦虑或愤怒等强烈情绪的内容往往更容易吸引关注，也可能因较高的互动量获得更多传播机会。当用户长时间浏览此类内容时，可能逐渐将其中反复强调的极端案例理解为普遍现象，进而高估现实风险，对自身处境产生不必要的担忧。

针对这一问题，破茧提供情绪引导识别与提醒功能。系统结合内容引用的事实及其适用条件，判断评价性结论是否存在范围外推或确定程度提高，并进一步识别其中的负面情绪强化方式。只有当相关推断越界与负面强化同时成立、且原文依据得到核验时，才将其标记为过度推断型情绪引导；随后结合浏览时间统计用户接触此类内容的情况。

例如，一名正在准备求职的学生，可能连续观看大量强调就业困难的视频，其中反复出现``毕业就失业''等绝对化判断。当系统识别到视频将个别招聘困难推断为普遍且不可改变的就业困境，并以灾难化表达强化担忧时，可以提示用户近期反复接触了此类内容，同时展示被引用的事实与相关结论，供用户核查。

用户还可以通过每日报告和长期分析报告，查看过度推断型情绪引导内容在浏览数量和有效浏览时长中的占比，了解相关表达主要集中于哪些主题。报告展示事实范围扩大、结论确定程度提高及负面情绪强化等具体特征，并保留对应原文。系统提供提醒设置，使用户能够根据自身需要调整提示方式和频率，避免影响正常浏览。

\textbf{（4）信息盲区提示：个人浏览记录中的观点覆盖局限提示}

仅仅识别信息茧房，并不能直接解决用户信息来源单一的问题。当系统发现某一主题长期集中于少数观点时，用户仍然需要自行寻找其他资料。然而，普通用户往往不知道应该从哪里开始查找，也难以判断哪些资料真正提供了不同的事实和解释。

为此，破茧提供面向具体议题的主动信息补全功能。当系统识别到某个主题的信息长期重复时，用户可以选择查看其他值得关注的观点，以及支持这些观点的不同事实依据。系统结合用户已经浏览的内容，检索尚未充分涉及的信息，并根据来源与内容之间的关系进行整理。

例如，用户近期接触的就业信息主要集中于互联网企业裁员。通过信息补全功能，系统可以进一步展示不同行业的招聘变化，补充相关地区的就业统计，并提供与已有观点不同的解释。对于引用相同新闻或统计资料的内容，系统也会提示其来源关系，避免将重复转载误认为相互独立的证据。

这一功能并不简单地向用户推荐与现有观点相反的内容，而是帮助其获得此前缺少的事实与解释。用户可以自主决定是否阅读相关资料，并通过原始来源核查其可靠性。

通过将信息茧房预警与主动信息补全结合起来，破茧能够从单纯发现信息重复，进一步延伸到帮助用户主动拓宽信息来源，使认知体检真正成为改善信息获取方式的工具。

\textbf{（5）主题观点整理：同主题观点归纳与论据来源追溯}

短视频平台以连续推荐和快速切换为主要内容呈现方式，用户每天可能浏览大量不同主题的信息。即使多次接触同一个问题，相关内容也往往分散在不同的浏览时段，中间穿插着大量无关视频。这种碎片化、非连续的信息获取方式，使用户难以将前后看到的观点联系起来，也很难长期记住不同内容提供的事实和论据。即使花费大量时间浏览，用户仍可能无法系统回顾自己对某个问题究竟获得了哪些认识，难以将零散的信息积累转化为有效的知识。

针对这一问题，破茧提供面向具体主题的观点归纳与论据整理功能。系统利用已有的浏览记录，将分散在不同时间和视频中的同主题内容汇总整理，并对相同或相近的观点进行归类。对于每种观点，系统统计支持该观点的已浏览内容数量，提取相关内容采用的主要论据，并保留对应的原始来源。同时，系统能够识别不同内容之间的论据重复关系，展示同一主题下存在的观点分歧，帮助用户了解各类观点的主要依据及其相互联系。

通过这一功能，用户无需重新翻看大量视频，就能够集中回顾某个主题下已经接触的观点及其依据，将原本分散、不连续的信息整理为结构清晰、便于回顾的主题总结。这不仅能够减少重复查找和阅读的时间，还能帮助用户建立不同信息之间的联系，加深对相关问题的理解。对于就业、升学等需要长期关注的现实问题，用户可以持续回顾和比较已有信息，发现尚需进一步查证的事实与论据，使日常碎片化浏览获得的信息逐步转化为可以积累、复用的知识，为后续学习和个人决策提供参考。

\textbf{（6）隐私保护与自主控制：浏览数据机密处理与用户授权管理}

个人浏览记录能够反映用户长期关注的问题，可能涉及多项个人生活中的敏感事项。对于需要持续分析浏览历史的工具而言，如果用户必须将全部记录直接交给服务提供者，就可能因为担心隐私泄露而不愿使用相关功能。

针对这一问题，破茧将可验证异构机密计算技术融入个人认知体检服务。客户端通过开源审计与版本验证约束采集程序行为，上传前利用远程证明核验云端可信程序及其CPU、GPU运行环境。浏览内容经加密提交后，只在通过验证的受保护程序内解密和处理，历史记录则以用户掌握密钥的密文形式存储。

在真实可信硬件及程序验证机制有效的条件下，包括作品运营团队在内的服务提供者，也无法仅凭普通服务器管理权限读取用户的明文浏览记录或长期分析结果。用户可以在保留历史数据解密控制权的同时，使用持续的认知体检服务。

通过上述功能，破茧形成了从个人信息体检到风险提醒，再到观点整理与信息盲区提示的完整服务过程。用户不仅能够了解每天看到了什么，还能够发现长期重复的信息，并将已经阅读的内容整理为清晰的观点总结。在此过程中，可验证机密计算技术为持续分析提供隐私保障，使用户能够更加放心地了解和管理自己的信息环境。

\subsubsection{作品特色}\label{section:1.4.2}

与传统数字健康工具和一般内容分析系统相比，破茧的特色不仅在于增加新的分析指标，更在于将信息茧房识别与个人信息管理结合起来，并通过可验证机密计算技术保护长期浏览数据。作品既关注用户日常接触的信息是否单一，也关注用户能否真正理解和利用这些分析结果。

\textbf{（1）从屏幕时间管理延伸到个人信息健康管理}

传统数字健康工具主要记录用户使用设备或应用的时间，能够帮助用户管理网络使用习惯，却难以说明这些时间具体带来了什么信息。破茧进一步分析用户实际浏览的内容，识别其中的观点及其论据，并结合浏览时间呈现个人信息环境的特点。

用户不再只能看到每天使用网络多长时间，还能够了解这些时间主要用于关注哪些问题，相关内容是否反复引用相同材料，以及是否反复接触借助局部事实夸大负面结论的内容。这使数字健康管理从关注设备使用行为，进一步延伸到关注个人实际获得的信息。

\textbf{（2）从被动的风险识别延伸到主动的认知安全防护}

普通内容分析工具通常在完成分类或风险判断后给出结果，用户仍需要自行理解问题并寻找改善方法。破茧则将长期信息茧房识别与主动提醒结合起来，使用户能够及时了解个人浏览记录中值得关注的变化。

在发现信息长期重复或过度推断型情绪引导反复出现时，系统不仅提供风险提示，还能够指出相关问题主要出现在哪些主题中。用户可以结合信息盲区提示，了解当前浏览内容可能存在哪些局限，并自主选择是否补充其他资料。

这一特点使破茧不再局限于对过去浏览情况的统计，而能够为用户调整信息获取方式提供持续参考。

\textbf{（3）将分散的信息整理为具有原始依据的观点总结}

网络内容通常以短视频和帖子等形式分散呈现，用户很难将大量零散信息整理成清晰的认识。破茧利用内容理解和历史比较结果，按照主题归纳用户已经阅读的观点，并统计每种观点获得了多少条已浏览内容的支持。

与普通的文章摘要不同，这种整理不仅概括单条内容，还能够将不同来源表达的相近观点合并，并保留相应的论据及原始来源。用户可以直接查看不同观点之间的分歧，也能够进一步判断支持同一观点的内容是否反复引用相同材料。

因此，破茧不仅帮助用户发现信息重复，还帮助用户真正记住和整理已经获得的信息，使认知体检具有更直接的日常使用价值。

\textbf{（4）将观点论据分析与情绪引导识别结合，呈现信息表达的潜在风险}

信息茧房不仅可能表现为相似观点和论据的反复出现，还可能伴随着对局部事实的绝对化解释和负面情绪强化。破茧复用信息茧房识别阶段提取的事实、观点及论证关系，判断相关结论是否超出原始材料的适用范围，以及其表达是否采用灾难化或无助化等方式放大负面情绪。

用户既可以了解某种观点是否被反复强调，也能够看到相关内容为何被判定为过度推断型情绪引导，以及具体的事实和结论原文。将观点分布、历史事实增量与情绪引导接触情况结合起来，作品能够从事实依据和表达方式两个角度展示个人信息环境。

\textbf{（5）面向高频信息流的轻量化快速处理能力}

短视频和社交平台的内容切换频繁，新浏览记录可能在很短时间内连续产生。如果每条内容都要求大模型生成完整分析，再与全部历史逐条比较，处理开销就会随着浏览量增加而不断扩大。破茧将连续分析拆分为片段定位、固定类别判别、历史候选比较和指标计算，减少重复的长文本生成与全量历史扫描。

在具体处理上，系统使用UIE-base快速定位观点与论据，由Simple Jev和SemIf完成本地结构化关系判断，并通过BM25与语义检索缩小历史比较范围。情绪引导模块复用已提取的事实及结论，通过共享模型完成多专家候选评分。对于重复出现的内容，系统复用已有的语义判别结果；不同新内容的抽取任务可以并行处理，涉及个人历史的更新则保持浏览顺序。由此在保留原文证据与长期比较能力的同时，提高对持续高频浏览记录的处理效率。

\textbf{（6）通过可验证异构机密计算减少对服务运营方的隐私信任依赖}

传统云端分析服务通常要求用户信任服务提供者不会查看或滥用其上传的数据。破茧则围绕采集、上传、计算和存储建立可验证的防护链路：通过客户端代码审计与版本验证约束采集行为，借助远程证明核验云端可信程序，并以TDX和机密GPU保护模型推理中的明文。

用户长期主密钥保留在手机端，历史索引与分析报告以密文存储；后续分析所需的数据密钥只在通过验证的会话中按次授权。在可信硬件与程序审核机制有效的条件下，普通服务运营权限不足以读取受保护的浏览原文及长期分析记录，使个人认知体检减少对服务提供者隐私承诺的依赖。

总体而言，破茧将信息茧房识别、情绪引导识别与观点整理统一于持续的认知体检服务之中，使用户能够发现信息环境可能存在的局限，并回顾已经获得的判断依据。作品通过轻量化的结构化决策与历史结果复用适应高频浏览需求，同时以可验证异构机密计算保护敏感数据，形成兼顾使用效率、日常实用性与认知安全防护的个人信息管理工具。

\subsection{典型应用场景与发展前景}\label{section:1.5}

\subsubsection{面向大众的日常信息健康管理场景}\label{section:1.5.1}

随着短视频和社交媒体深度融入日常生活，普通用户每天都会接触大量网络内容。然而，这些信息通常分散在不同平台和时间段，用户很难完整回忆自己看到了什么，更难判断哪些内容只是对已有信息的重复。对于普通用户而言，网络浏览管理的需求已经不再局限于减少屏幕使用时间。人们还希望了解自己每天获得了哪些有价值的信息，是否长时间沉浸在相似的内容中，以及已经阅读的材料能否得到有效整理。

针对这一需求，破茧将个人认知体检融入日常浏览过程。用户在授权后，可以通过24小时信息分布图了解当天主要关注的主题，并进一步查看不同观点的分布和信息重复情况。对于反复出现的过度推断型情绪引导内容，系统提供相应提示，帮助用户认识日常信息环境的特点。

例如，一名正在寻找工作的大学生，可能每天花费大量时间浏览招聘信息和求职经验。当晚的认知体检报告能够展示就业相关内容所占的浏览时间，以及其中主要出现了哪些观点。用户还可以打开主题观点总结，了解当天看到的不同就业观点分别由多少条内容支持，并查看相关论据。如果大量视频只是重复讨论同一则裁员新闻，系统会标注相关内容之间的重复关系。对于将局部困难推断为普遍危机并强化担忧的内容，系统也能够展示其出现比例与累计浏览时长，使用户更加清楚地了解自己当天主要接触了什么样的信息。

这一场景的价值在于，用户不必重新翻看大量视频，就能够回顾当天获得的信息，并根据需要主动补充其他资料。破茧由此成为个人日常信息管理的辅助工具，使网络浏览从单纯的内容消费，进一步转化为可以回顾和整理的信息积累过程。同时，作品通过可验证机密计算保护个人浏览记录，使用户能够在减少向服务运营方暴露明文信息的条件下使用相关服务，为面向大众提供持续的个人认知体检建立必要的隐私保护基础。

\subsubsection{面向重大公共议题的个人认知安全防护场景}\label{section:1.5.2}

在重大公共事件和社会热点发生期间，网络信息通常具有传播速度快、观点分歧明显的特点。特别是在突发自然灾害和公共卫生事件等情境下，用户可能在短时间内接触大量相关内容，其中既包括事实报道和专业分析，也可能包含未经核实的说法及强烈的情绪表达。由于此类事件通常与公众切身利益密切相关，容易激发担忧或愤怒等情绪，相关内容也可能因较高的互动量获得更多传播机会。当不同账号反复引用同一事件，却缺少新的事实和独立来源时，用户可能误以为某种判断得到了广泛证实。如果再叠加持续出现的情绪化表达，便可能使用户高估局部风险，忽视事件的其他发展情况，进而影响其对公共事件的理解和判断。

相比日常兴趣内容，重大公共议题中的信息失衡可能对现实决策产生更加直接的影响。例如，在突发灾害期间，用户如果长期看到未经充分核实的极端案例，却很少接触权威部门发布的最新信息，就可能对实际风险形成片面认识，甚至作出不必要或不恰当的应对决定。

针对这一场景，破茧能够对用户在特定公共议题下的浏览记录进行持续分析。系统识别相关内容的观点及其论据，并比较不同账号是否反复引用相同材料。同时，结合浏览时间和情绪引导识别结果，帮助用户了解自己是否长期接触某一种解释，以及相关结论是否持续通过不合理外推强化负面情绪。当某一公共事件持续引发社会关注时，用户还可以通过主题观点整理功能，查看自己当天围绕该事件接触了哪些不同观点，以及每种观点获得了多少条已浏览内容的支持。对于明显存在争议的内容，系统保留其对应的原始论据，使用户能够进一步核查相关事实。

需要强调的是，破茧不直接根据观点出现的次数判断新闻真实性，也不会将情绪强烈的内容自动认定为恶意信息。作品主要分析个人已经接触的信息是否过于单一，以及不同内容之间是否存在重复传播。用户可以结合报告主动查阅权威信息或其他独立来源，减少仅凭反复出现的内容作出判断的风险。通过帮助用户识别重复信息和过度推断型情绪引导，破茧能够将日常信息管理延伸到重大公共议题下的个人认知安全防护，为用户辨别网络信息提供辅助支持。

\subsubsection{面向网络认知安全教育的信息素养提升场景}\label{section:1.5.3}

随着网络内容生产和传播方式日益复杂，公众不仅需要掌握搜索和获取信息的能力，还需要学会判断不同信息是否具有独立依据，并认识情绪化表达可能对个人判断产生的影响。传统的信息素养教育通常通过案例讲解和知识培训，向公众介绍信息茧房与算法推荐的基本原理。然而，用户即使了解这些概念，也未必能够判断自己日常使用网络时是否遇到了类似情况。缺少与个人实际浏览经历的联系，容易使相关教育停留在理论层面，难以帮助用户形成持续反思信息来源的习惯。

针对这一问题，破茧可以为高校网络安全教育和数字素养培训提供辅助工具。通过用户自愿参与的认知体检活动，系统将其实际浏览的信息整理为可视化报告，使用户能够直接看到自己关注的主题，以及不同观点在浏览内容中的分布情况。例如，在高校信息素养教育活动中，参与者可以自主查看个人认知体检报告，了解近期是否大量接触相似观点，以及不同内容是否反复引用同一事件。借助主题观点整理功能，参与者还可以观察同一个问题下存在怎样的观点分歧，并尝试辨别浏览内容数量与独立论据数量之间的区别。

在后续使用中，参与者可以自主调整信息获取方式，并对比前后报告中观点分布与信息重复情况的变化。通过这种方式，用户能够在真实浏览经历中理解信息茧房的形成特点，逐步培养主动查证信息和关注不同事实依据的意识。

整个过程以个人自愿参与和自主查看结果为原则，不需要向教师或组织管理者开放原始浏览历史。破茧通过隐私保护机制减少第三方接触个人数据的机会，使认知安全教育能够在尊重个人隐私的基础上开展。这一应用场景将个人认知体检与信息素养教育结合起来，使抽象的认知安全问题转化为可以观察和讨论的实际案例，为公众数字素养和个人认知安全防护能力的提升提供技术支持。

\subsubsection{应用价值与发展前景}\label{section:1.5.4}

上述应用场景体现了破茧在个人生活和认知安全领域的不同价值。对于普通互联网用户，作品能够帮助其整理零散阅读内容，减少重复浏览带来的时间消耗，并认识个人信息环境中的潜在局限。对于重大公共议题，作品能够辅助用户了解不同观点的分布及其论据，减少因相同信息反复传播而产生的误判。对于认知安全教育，作品则提供了将理论知识与个人真实浏览经历相结合的实践工具。

从用户需求来看，破茧兼顾认知安全防护与日常信息管理。用户即使没有明确察觉到信息茧房风险，也可以通过主题观点整理功能获得实际帮助。对于需要长期关注某个问题的用户，系统能够将分散在不同时间的阅读内容整理为可回顾的信息记录，从而提高网络信息的利用效率。

从技术应用来看，作品充分利用已经获得的浏览记录及其分析结果，实现信息茧房识别与主题观点整理等功能。信息盲区提示同样以个人历史记录为基础，不依赖持续的外部网络检索。这种方式能够减少额外的数据获取和模型推理开销，为后续在个人设备和云端服务之间合理分配计算任务提供条件。

从信息安全角度来看，破茧通过可验证机密计算机制保护个人敏感数据，减少服务运营方直接接触明文浏览记录的机会。随着可信计算基础设施和人工智能服务技术的发展，这种在隐私保护条件下开展个人长期信息分析的方式，也具有进一步推广的应用潜力。

总体而言，破茧探索了一种以用户为中心的认知安全防护方式。作品不要求用户放弃个性化推荐服务，也不替代用户对信息真伪和价值的最终判断，而是通过持续分析和清晰的信息整理，帮助用户发现自己可能忽略的观点，理解不同内容之间的联系，并更加主动地选择信息来源。

\textbf{通过将信息茧房识别、情绪引导识别与提醒、主题观点整理及可验证异构机密计算相结合，破茧既能够满足普通用户日常信息管理的需求，也能够为网络空间中的个人认知安全防护提供技术支持，使认知安全从抽象的研究问题转化为面向个人的实际服务能力。}
